\documentclass[10pt,a4paper]{amsart}
\usepackage{amsmath,amssymb,amsthm,graphicx}
% Retained preamble for the Stacks Project chapter {CH} (shared project preamble).
% The chapter file pinned-original-source.tex only carries \input{{preamble}} on its first line,
% so no preamble text lives inside the chapter itself.  The Stacks Project uses one single shared
% preamble for every chapter: the snapshot stages it three times (sources/stacks-chow/preamble.tex,
% sources/stacks-etale-cohomology/preamble.tex and focused-seeds/AGP/sources/stacks-snapshot-20261011/
% preamble.tex) and those copies are identical in size.  The block below is copied mechanically,
% line for line, from sources/stacks-chow/preamble.tex lines 14-18 and 154-217: the Xy-pic setup and
% the native theorem declarations together with the mathematical macros.  Nothing was added, removed
% or redefined; the wrapper class and the amsmath/amssymb/amsthm/graphicx base packages are supplied
% by the assembler.
% For commutative diagrams we use Xy-pic
% Chapter: stacks-spaces-more-morphisms.tex -- More on Morphisms of Spaces.
\usepackage[all]{xy}

% We use 2cell for 2-commutative diagrams.
\xyoption{2cell}

% Theorem environments.
%
\theoremstyle{plain}
\newtheorem{theorem}[subsection]{Theorem}
\newtheorem{proposition}[subsection]{Proposition}
\newtheorem{lemma}[subsection]{Lemma}

\theoremstyle{definition}
\newtheorem{definition}[subsection]{Definition}
\newtheorem{example}[subsection]{Example}
\newtheorem{exercise}[subsection]{Exercise}
\newtheorem{situation}[subsection]{Situation}

\theoremstyle{remark}
\newtheorem{remark}[subsection]{Remark}
\newtheorem{remarks}[subsection]{Remarks}

\numberwithin{equation}{subsection}

% Macros
%
\def\lim{\mathop{\mathrm{lim}}\nolimits}
\def\colim{\mathop{\mathrm{colim}}\nolimits}
\def\Spec{\mathop{\mathrm{Spec}}}
\def\Hom{\mathop{\mathrm{Hom}}\nolimits}
\def\Ext{\mathop{\mathrm{Ext}}\nolimits}
\def\SheafHom{\mathop{\mathcal{H}\!\mathit{om}}\nolimits}
\def\SheafExt{\mathop{\mathcal{E}\!\mathit{xt}}\nolimits}
\def\Sch{\mathit{Sch}}
\def\Mor{\mathop{\mathrm{Mor}}\nolimits}
\def\Ob{\mathop{\mathrm{Ob}}\nolimits}
\def\Sh{\mathop{\mathit{Sh}}\nolimits}
\def\NL{\mathop{N\!L}\nolimits}
\def\CH{\mathop{\mathrm{CH}}\nolimits}
\def\proetale{{pro\text{-}\acute{e}tale}}
\def\etale{{\acute{e}tale}}
\def\QCoh{\mathit{QCoh}}
\def\Ker{\mathop{\mathrm{Ker}}}
\def\Im{\mathop{\mathrm{Im}}}
\def\Coker{\mathop{\mathrm{Coker}}}
\def\Coim{\mathop{\mathrm{Coim}}}

% Boxtimes
%
\DeclareMathSymbol{\boxtimes}{\mathbin}{AMSa}{"02}

%
% Macros for moduli stacks/spaces
%
\def\QCohstack{\mathcal{QC}\!\mathit{oh}}
\def\Cohstack{\mathcal{C}\!\mathit{oh}}
\def\Spacesstack{\mathcal{S}\!\mathit{paces}}
\def\Quotfunctor{\mathrm{Quot}}
\def\Hilbfunctor{\mathrm{Hilb}}
\def\Curvesstack{\mathcal{C}\!\mathit{urves}}
\def\Polarizedstack{\mathcal{P}\!\mathit{olarized}}
\def\Complexesstack{\mathcal{C}\!\mathit{omplexes}}
% \Pic is the operator that assigns to X its picard group, usage \Pic(X)
% \Picardstack_{X/B} denotes the Picard stack of X over B
% \Picardfunctor_{X/B} denotes the Picard functor of X over B
\def\Pic{\mathop{\mathrm{Pic}}\nolimits}
\def\Picardstack{\mathcal{P}\!\mathit{ic}}
\def\Picardfunctor{\mathrm{Pic}}
\def\Deformationcategory{\mathcal{D}\!\mathit{ef}}
\newenvironment{reference}{}{}
\newenvironment{slogan}{}{}
\newenvironment{history}{}{}
\usepackage{cleveref}
\crefname{originallemma}{original lemma}{original lemmas}
\Crefname{originallemma}{Original lemma}{Original lemmas}
\makeatletter
\newlabel{lemma-base-change-fibres-nr-geometrically-connected-components}{{055F}{source}{}{sourcepointer-1}{}}
\newlabel{lemma-base-change-fibres-nr-geometrically-connected-components@cref}{{[originallemma][0][]055F}{[1][0][]source}{}{}{}}
\makeatother
\crefname{originallemma}{original lemma}{original lemmas}
\Crefname{originallemma}{Original lemma}{Original lemmas}
\makeatletter
\newlabel{lemma-stein-universally-closed-residue-fields}{{0E0M}{source}{}{sourcepointer-2}{}}
\newlabel{lemma-stein-universally-closed-residue-fields@cref}{{[originallemma][0][]0E0M}{[1][0][]source}{}{}{}}
\makeatother
\crefname{originallemma}{original lemma}{original lemmas}
\Crefname{originallemma}{Original lemma}{Original lemmas}
\makeatletter
\newlabel{more-morphisms-lemma-etale-makes-integral-split}{{0BSR}{source}{}{sourcepointer-3}{}}
\newlabel{more-morphisms-lemma-etale-makes-integral-split@cref}{{[originallemma][0][]0BSR}{[1][0][]source}{}{}{}}
\makeatother
\crefname{originallemma}{original lemma}{original lemmas}
\Crefname{originallemma}{Original lemma}{Original lemmas}
\makeatletter
\newlabel{spaces-morphisms-lemma-fppf-open}{{042S}{source}{}{sourcepointer-4}{}}
\newlabel{spaces-morphisms-lemma-fppf-open@cref}{{[originallemma][0][]042S}{[1][0][]source}{}{}{}}
\makeatother
\crefname{originallemma}{original lemma}{original lemmas}
\Crefname{originallemma}{Original lemma}{Original lemmas}
\makeatletter
\newlabel{theorem-stein-factorization-general}{{03H2}{source}{}{sourcepointer-5}{}}
\newlabel{theorem-stein-factorization-general@cref}{{[originallemma][0][]03H2}{[1][0][]source}{}{}{}}
\makeatother
\title{The number of geometric connected components of fibres of a proper flat finitely presented map is lower semi-continuous}
\author{The Stacks Project Authors}
\date{stacks-project:spaces-more-morphisms.tex:local-snapshot-20261011; GNU FDL 1.2 or later}
\begin{document}
\maketitle
\noindent\textit{Editorial note.} Statement span is lines 8349-8358 of sources/stacks-spaces-more-morphisms/pinned-original-source.tex; proof span is lines 8360-8384, the complete author proof (23 body lines).\par\medskip
\noindent\textit{Editorial note.} Navigation stubs: 5 cited-but-not-included labels are supplied as non-load-bearing pointers with their pinned tags where the snapshot tag index has them.\par\medskip
\noindent\textit{Editorial note.} extra\_packages is empty on purpose: the retained preamble already loads xy with options and a second declaration clashes in the assembler wrapper.\par\medskip
\noindent\textit{Editorial note.} Text is retained verbatim, including the author's own typos and placeholders.\par\medskip
\noindent\textit{Editorial note.} Counted B side (both complete); depends on L16200 (same wave) and one queue card (wave 21).\par\medskip
\noindent\textit{Allowed standard prerequisites (editorial).}\begin{itemize}
\item Explicit prerequisite boundary. Chapter theorems/lemmas stein-factorization-general (L16200), stein-universally-closed-residue-fields, base-change-fibres-nr-geometrically-connected-components; More on Morphisms, lemma etale-makes-integral-split; Morphisms of Spaces, lemma fppf-open.
\end{itemize}\medskip
\ifdefined\hypertarget\hypertarget{sourcepointer-1}{}\fi
\noindent\textit{Source reference.} Original lemma 055F: sources/stacks-spaces-more-morphisms/pinned-original-source.tex, source lines 7070--7084. Cited by the retained statement or proof but not part of the collected spans; supplied only as a navigation pointer. The argument does not depend on its content.\par\medskip
\ifdefined\hypertarget\hypertarget{sourcepointer-2}{}\fi
\noindent\textit{Source reference.} Original lemma 0E0M: sources/stacks-spaces-more-morphisms/pinned-original-source.tex, source lines 8061--8066. Cited by the retained statement or proof but not part of the collected spans; supplied only as a navigation pointer. The argument does not depend on its content.\par\medskip
\ifdefined\hypertarget\hypertarget{sourcepointer-3}{}\fi
\noindent\textit{Source reference.} Original lemma 0BSR: sources/stacks-spaces-more-morphisms/pinned-original-source.tex, source lines 8384--8384. Cited by the retained statement or proof but not part of the collected spans; supplied only as a navigation pointer. The argument does not depend on its content.\par\medskip
\ifdefined\hypertarget\hypertarget{sourcepointer-4}{}\fi
\noindent\textit{Source reference.} Original lemma 042S: sources/stacks-spaces-more-morphisms/pinned-original-source.tex, source lines 8384--8384. Cited by the retained statement or proof but not part of the collected spans; supplied only as a navigation pointer. The argument does not depend on its content.\par\medskip
\ifdefined\hypertarget\hypertarget{sourcepointer-5}{}\fi
\noindent\textit{Source reference.} Original lemma 03H2: sources/stacks-spaces-more-morphisms/pinned-original-source.tex, source lines 8197--8214. Cited by the retained statement or proof but not part of the collected spans; supplied only as a navigation pointer. The argument does not depend on its content.\par\medskip
\begin{lemma}
\label{lemma-proper-flat-nr-geom-conn-comps-lower-semicontinuous}
Let $S$ be a scheme.
Let $X \to Y$ be a morphism of algebraic spaces over $S$.
If $f$ is proper, flat, and of finite presentation, then the function
$n_{X/Y} : |Y| \to \mathbf{Z}$ counting the number of geometric
connected components of fibres of $f$
(Lemma \ref{lemma-base-change-fibres-nr-geometrically-connected-components})
is lower semi-continuous.
\end{lemma}

\begin{proof}
The question is \'etale local on $Y$, hence we may and do assume $Y$
is an affine scheme. Let $y \in Y$. Set $n = n_{X/S}(y)$.
Note that $n < \infty$ as the geometric fibre of $X \to Y$ at $y$
is a proper algebraic space over a field, hence Noetherian, hence
has a finite number of connected components.
We have to find an open neighbourhood $V$ of $y$ such that $n_{X/S}|_V \geq n$.
Let $X \to Y' \to Y$ be the Stein factorization as in
Theorem \ref{theorem-stein-factorization-general}.
By Lemma \ref{lemma-stein-universally-closed-residue-fields}
there are finitely many points
$y'_1, \ldots, y'_m \in Y'$ lying over $y$
and the extensions $\kappa(y'_i)/\kappa(y)$ are finite.
More on Morphisms, Lemma \ref{more-morphisms-lemma-etale-makes-integral-split}
tells us that after replacing $Y$ by an \'etale neighbourhood
of $y$ we may assume $Y' = V_1 \amalg \ldots \amalg V_m$ as a scheme
with $y'_i \in V_i$ and $\kappa(y'_i)/\kappa(y)$ purely inseparable.
Then the algebraic spaces $X_{y_i'}$ are geometrically connected over
$\kappa(y)$, hence $m = n$. The algebraic spaces
$X_i = (f')^{-1}(V_i)$, $i = 1, \ldots, n$
are flat and of finite presentation over $Y$. Hence the image of $X_i \to Y$
is open (Morphisms of Spaces, Lemma \ref{spaces-morphisms-lemma-fppf-open}).
Thus in a neighbourhood of $y$ we see that $n_{X/Y}$ is
at least $n$.
\end{proof}

\end{document}
